\section{Reviewer Navigation Matrix and Audit Checklist}
\label{sec:oc13-reviewer-navigation}

Start with the row that matches the question or concern, inspect the primary
route, and use the verification column as the pass/fail test.
This appendix uses three canonical artifact terms:
\begin{itemize}[leftmargin=1.6em,nosep]
\item \emph{Master Monograph}: the English flagship volume. Later uses of
      \emph{Monograph} are shorthand for this same artifact.
\item \emph{Journal Core}: the shorter bounded extraction from the monograph.
\item \emph{Release}: the outward package that also contains supporting
      surfaces, PDFs, assets, and companion dossiers.
\end{itemize}
Capitalisation marks the canonical glossary terms above; lowercase uses are
generic prose shorthand unless explicitly styled as one of those terms.
The following class labels are local definitions for this appendix rather than
additional canonical glossary terms.
Concrete artifact and witness classes are as follows:
\begin{itemize}[leftmargin=1.6em,nosep]
\item \emph{Surfaces} are the release-bundled
      \occode{editorial/*_latest.json} projections.
\item \emph{Companion dossiers} are the generated dossier packages under
      \occode{editorial/dossier_packages/}.
\item \emph{Assets} are the released figures and tables under
      \occode{assets/}.
\item \emph{Archival PDF witness}: the prior Core~1.2 PDF named in the
      journal-core provenance route in
      Appendix~\ref{sec:oc13-journal-core-bridge}: \emph{Ontology of
      Continua --- Core v1.2.0}.
\end{itemize}

\subsection{Review questions, compact route, and decisive checks}

\begingroup
\footnotesize
\RaggedRight
\setlength{\emergencystretch}{1.5em}
\setlength{\tabcolsep}{3pt}
\setlength{\LTleft}{0pt}
\setlength{\LTright}{0pt}

\paragraph{Core theorem route and scope.}
\par{\centering
\begin{tabular}{@{}L{0.18\textwidth}L{0.26\textwidth}L{0.46\textwidth}@{}}
\toprule
\textbf{Reviewer question} & \textbf{Primary route to inspect} & \textbf{Verification cue} \\
\midrule
What is the central defended result? & Abstract; Sections~\ref{sec:results},~\ref{sec:oc13-theorem-roadmap} & Theorem~A: bounded post-collapse identity classification; no general lift theorem. \\
Where does the proof-bearing route begin? & Sections~\ref{sec:results},~\ref{sec:oc13-theorem-roadmap} & Verify that Axiom~3.2 with Definitions~12.1, 12.3, and 12.4 supports Source Theorem~3; that Source Theorem~3 supports Lemma~1; that Lemma~1 with Definition~12.5 supports Lemma~2; and that Lemma~2 supports Theorem~A. Inspect Axiom~3.3, Definition~12.6, and Source Corollary~3.2 separately as the companion rebirth-control branch. \\
Is collapse defined formally or rhetorically? & Sections~\ref{sec:results},~\ref{sec:collapse-rebirth} & Admissible-state emptiness, \(k\to0\), and threshold failure. \\
What exactly is residue? & Sections~\ref{sec:collapse-rebirth},~\ref{sec:oc13-worked-examples} & Post-collapse structural trace; no continuing live identity. \\
What exactly is rebirth? & Sections~\ref{sec:collapse-rebirth},~\ref{sec:oc13-worked-examples} & Later live continuum grounded in residue; numerically new. \\
Why is irreversibility necessary? & Sections~\ref{sec:results},~\ref{sec:oc13-theorem-roadmap} & Blocks same-entity restoration after lawful death. \\
Could the main theorem depend on an excluded lift-heavy theorem family? & Section~\ref{sec:oc13-theorem-roadmap}; Appendix~\ref{sec:oc13-journal-core-bridge} & No excluded lift-heavy theorem family is used in the bounded proof route. \\
Why is the monograph larger than the journal core? & Section~\ref{sec:oc13-publication-boundary}; Appendix~\ref{sec:oc13-journal-core-bridge} & Didactic shell, formal architecture, and repair logic live here. \\
Why is the journal core smaller without becoming weaker? & Appendix~\ref{sec:oc13-journal-core-bridge} & One-way extraction; theorem scope unchanged. \\
Why is the hierarchy fixed as \(K_0,\dots,K_{12}\)? & Sections~\ref{sec:oc13-foundational-consistency},~\ref{sec:klevels-full},~\ref{sec:oc13-worked-examples} & Upper levels carry substantive higher-order roles. \\
\bottomrule
\end{tabular}
\par}


\subsection{External criticism evidence route}
\label{sec:oc13-reviewer-navigation-external-criticism-v3}

\begin{longtable}{@{}p{0.20\textwidth}p{0.28\textwidth}p{0.42\textwidth}@{}}
\toprule
Reviewer question & Primary route to inspect & Verification cue \\
\midrule
\endfirsthead
\toprule
Reviewer question & Primary route to inspect & Verification cue \\
\midrule
\endhead
Did external criticism change the release? & Appendix~\ref{sec:oc13-external-criticism-closure}; Section~\ref{sec:oc13-source-tradition-remediation-v3} & Pass only if the response names concrete manuscript, bibliography, evidence, replay, theorem-fate, tuple/K rigor, or release-packet deltas. \\
Is Stanislav Tsukrov endorsed as approving the work? & Public acknowledgement packet and owner-gated no-send boundary & No. The acknowledgement records substantive review pressure, not endorsement or publication approval. \\
Can predictive claims be read strongly? & Section~\ref{sec:oc13-prediction-claim-taxonomy-v3}; Appendix~\ref{sec:oc13-external-criticism-closure} & No quantitative prediction without protocol, data/code, output, comparator, and falsifier. \\
\bottomrule
\end{longtable}

\clearpage
\paragraph{Audit, provenance, and review posture.}
\par{\centering
\begin{tabular}{@{}L{0.18\textwidth}L{0.26\textwidth}L{0.46\textwidth}@{}}
\toprule
\textbf{Reviewer question} & \textbf{Primary route to inspect} & \textbf{Verification cue} \\
\midrule
What would falsify the upper-level doctrine? & Sections~\ref{sec:oc13-worked-examples},~\ref{sec:klevels-full} & Lossless reduction of \(K_{11}\) and \(K_{12}\) to \(K_{10}\). \\
How do proof-oriented statement classes differ? & Sections~\ref{sec:oc13-foundational-consistency},~\ref{sec:oc13-theorem-roadmap} & Axiomatic, derived, structural, interpretive, empirical, and frontier rows keep separate burdens; editorial-bridge language remains packaging support rather than a proof class. \\
What is merely interpretive? & Sections~\ref{sec:discussion},~\ref{sec:oc13-publication-boundary},~\ref{sec:oc13-worked-examples} & Pass only if interpretive language stays outside the proof route and is not promoted into theorem, replay, or practical-use authority. \\
What is the exact source base? & Section~\ref{sec:oc13-publication-boundary}; Appendix~\ref{sec:oc13-source-audit-appendix} & Source corpus, SPOT surface, and archival PDF witness are separate objects. \\
Why is ORCID foregrounded? & Title page; Section~\ref{sec:oc13-publication-boundary} & Stable scholarly identity. \\
What is the role of companion dossiers? & Section~\ref{sec:oc13-publication-boundary}; Appendix~\ref{sec:oc13-journal-core-bridge} & Support and provenance exposure, not replacement authority. \\
How should a non-specialist enter the monograph? & Sections~\ref{sec:oc13-reader-contract},~\ref{sec:oc13-worked-examples} & Didactic ladder before proof audit. \\
Where are the most likely points of critique? & Section~\ref{sec:oc13-theorem-roadmap}; Appendix~\ref{sec:oc13-journal-core-bridge} & Collapse semantics, irreversibility, upper levels, extraction discipline. \\
What would count as a decisive release failure? & Sections~\ref{sec:oc13-worked-examples},~\ref{sec:falsifiability-extended},~\ref{sec:oc13-publication-boundary} & Continued identity after admissible-state loss, or reliance on an excluded theorem family. \\
What remains outside scope? & Sections~\ref{sec:discussion},~\ref{sec:oc13-publication-boundary}; Appendix~\ref{sec:oc13-journal-core-bridge} & Pass only if excluded theorem families, frontier rows, and projection targets remain outside the positive proof body. \\
\bottomrule
\end{tabular}
\par}

\endgroup

\subsection{Compact reading order by reviewer role}

\paragraph{Rapid critical reviewer.}
Read the abstract, theorem roadmap, relevant theorems and corollaries, the worked examples,
and the journal-core bridge appendix. Return to this matrix only to choose a
new route after that first pass. This route should reveal immediately
whether the release is too large, too narrow, or too broad in scope.

\paragraph{Formal mathematical reviewer.}
Read the model, results, axioms, operators, collapse/rebirth material, hierarchy chapter,
notation appendix, and source-audit chapter. The decisive question is whether the derived bounded
claim can be reconstructed from the stated source-native rows without informal supplementation.

\paragraph{Generalist scientific reviewer.}
Read the introduction, reader guide, theorem roadmap, worked examples, discussion, conclusion,
and then return to the formal chapters selectively. The decisive question here is whether the
volume explains itself honestly enough that scientific judgment can precede local notation
mastery.

\clearpage
